\section{Prediction of Oncogene Mutation Status in Non-Small Cell Lung Cancer: A Systematic Review and Meta-Analysis with a Special Focus on Artificial Intelligence-Based Methods}

Fuster-Matanzo et al. (2025) conducted a comprehensive systematic review and meta-analysis to evaluate artificial intelligence (AI)-based imaging approaches for predicting oncogene mutation status in non-small cell lung cancer (NSCLC)~\cite{Fuster-Matanzo2025-ki}. The study aimed to assess the potential of radiomics and deep learning models as non-invasive alternatives to conventional tissue-based molecular testing, focusing on key driver mutations including EGFR, ALK, and KRAS. 

The review followed PRISMA guidelines and included studies published up to June 30, 2023. Systematic searches were performed across MEDLINE (via PubMed), Embase, and the Cochrane Library. Studies were included if they used radiomics or AI-based methods to predict oncogene mutation status in NSCLC. From several hundred initially retrieved publications, 124 met inclusion criteria for the systematic review, and 51 were included in the quantitative meta-analysis. Analyses were stratified by mutation type and by model design-radiomics-only or radiomics combined with clinical data. Pooled sensitivity and false positive rate (FPR) were calculated using random-effects models, and meta-regression was applied to explore the influence of factors such as imaging modality, sample size, feature selection strategy, and inclusion of clinical variables.

For radiomics-only models based on CT-derived features, pooled performance metrics were as follows: for EGFR, sensitivity = 0.754 and FPR = 0.344; for ALK, sensitivity = 0.754 and FPR = 0.225; and for KRAS, sensitivity = 0.475 and FPR = 0.181. When radiomic and clinical features were combined (available for EGFR only), pooled sensitivity increased to 0.806 with a corresponding FPR of 0.315. The meta-regression did not identify any significant moderating effect of methodological or clinical factors on model performance.

The findings demonstrate that AI-driven radiomics models can achieve strong predictive performance for EGFR and ALK mutation status in NSCLC, with lower accuracy for KRAS. The modest improvement observed when incorporating clinical data suggests that imaging features alone carry substantial discriminative power. The authors highlighted that such non-invasive image-based biomarkers could complement or, in some cases, reduce the need for invasive biopsy procedures in clinical practice.

Nonetheless, several challenges remain before clinical translation, including heterogeneity in imaging acquisition and feature extraction, lack of standardization across studies, limited external validation, and variability in AI model reporting. The authors emphasized that future research should focus on multi-center validation, harmonization of imaging workflows, transparency in model development, and integration of explainable AI techniques.

Overall, this systematic review and meta-analysis by Fuster-Matanzo et al. underscored the growing role of AI-based radiomics in precision oncology, showing that CT-derived imaging signatures can non-invasively predict oncogene mutation status in NSCLC. However, achieving clinical applicability will require standardized methodologies, larger datasets, and consistent validation frameworks.